% !TeX encoding = UTF-8
% 吉田第六版：中译本第 9 章；此文件当前仅为目录与正文占位。
\SetChapterInfo{subtitle={吉田耕作第六版主线 · 正文待撰写}}
\chapter{半群的分析理论}\label{ch:yosida-09}
\YosidaChapterPlaceholder{9}{Analytical Theory of Semi-groups}{243--286}{231--273}

\section{\texorpdfstring{$(C_0)$}{(C0)} 类半群}\label{sec:yosida-09-01}
\YosidaSectionPlaceholder{243}{232}
% TODO: 定义与记号；定理与完整证明；例子；习题；中英文核对。

\section{局部凸空间中 \texorpdfstring{$(C_0)$}{(C0)} 类等度连续半群、几例半群}\label{sec:yosida-09-02}
\YosidaSectionPlaceholder{246}{234}
% TODO: 定义与记号；定理与完整证明；例子；习题；中英文核对。

\section{\texorpdfstring{$(C_0)$}{(C0)} 类等度连续半群的无穷小生成元}\label{sec:yosida-09-03}
\YosidaSectionPlaceholder{248}{237}
% TODO: 定义与记号；定理与完整证明；例子；习题；中英文核对。

\section{无穷小生成元 \texorpdfstring{$A$}{A} 的预解式}\label{sec:yosida-09-04}
\YosidaSectionPlaceholder{252}{240}
% TODO: 定义与记号；定理与完整证明；例子；习题；中英文核对。

\section{几例无穷小生成元}\label{sec:yosida-09-05}
\YosidaSectionPlaceholder{254}{242}
% TODO: 定义与记号；定理与完整证明；例子；习题；中英文核对。

\section{具等度连续幂的连续线性算子的指数函数}\label{sec:yosida-09-06}
\YosidaSectionPlaceholder{256}{244}
% TODO: 定义与记号；定理与完整证明；例子；习题；中英文核对。

\section{\texorpdfstring{$(C_0)$}{(C0)} 类等度连续半群用其无穷小生成元的表示和刻画}\label{sec:yosida-09-07}
\YosidaSectionPlaceholder{258}{246}
% TODO: 定义与记号；定理与完整证明；例子；习题；中英文核对。

\section{压缩半群和耗散算子}\label{sec:yosida-09-08}
\YosidaSectionPlaceholder{261}{250}
% TODO: 定义与记号；定理与完整证明；例子；习题；中英文核对。

\section{\texorpdfstring{$(C_0)$}{(C0)} 类等度连续群、Stone 定理}\label{sec:yosida-09-09}
\YosidaSectionPlaceholder{263}{251}
% TODO: 定义与记号；定理与完整证明；例子；习题；中英文核对。

\section{全纯半群}\label{sec:yosida-09-10}
\YosidaSectionPlaceholder{266}{254}
% TODO: 定义与记号；定理与完整证明；例子；习题；中英文核对。

\section{闭算子的分数幂}\label{sec:yosida-09-11}
\YosidaSectionPlaceholder{271}{259}
% TODO: 定义与记号；定理与完整证明；例子；习题；中英文核对。

\section{半群的收敛性、Trotter--Kato 定理}\label{sec:yosida-09-12}
\YosidaSectionPlaceholder{281}{269}
% TODO: 定义与记号；定理与完整证明；例子；习题；中英文核对。

\section{对偶半群、Phillips 定理}\label{sec:yosida-09-13}
\YosidaSectionPlaceholder{284}{272}
% TODO: 定义与记号；定理与完整证明；例子；习题；中英文核对。
